\chapter{Fix 11: Vacuum Uniqueness via Cluster Decomposition}
\label{ch:fix11_vacuum_uniqueness}

\section{Context: The Trap of Degenerate Vacua}
A Mass Gap $\Delta > 0$ is defined as the energy difference between the ground state (vacuum) $\Omega$ and the first excited state. However, if the vacuum is degenerate (e.g., due to Spontaneous Symmetry Breaking), the "gap" might be zero in the thermodynamic limit. To rigorously define the gap, we must prove the vacuum $\Omega$ is unique. This relies on \textbf{Appendix AK}.

\section{Theorem G.1 (Exponential Clustering)}
For the lattice Yang-Mills measure $d\mu$ at any coupling $\beta$ where the cluster expansion converges (or center symmetry is unbroken), the truncated two-point function satisfies:
\begin{equation}
|\langle \mathcal{O}_A(x) \mathcal{O}_B(y) \rangle - \langle \mathcal{O}_A(x) \rangle \langle \mathcal{O}_B(y) \rangle| \le C e^{-m |x-y|}
\end{equation}
where $m \ge \text{Re}(\Delta)$ is the mass gap. This implies the ground state of the infinite volume Hamiltonian is non-degenerate (a pure phase).

\section{Proof Derivation}
\subsection{1. Contour Representation}
We map the lattice gauge theory to a statistical mechanics model of "contours" (geometric defects) $\gamma$. The partition function becomes $Z = \sum_{\Gamma} \prod_{\gamma \in \Gamma} \rho(\gamma)$.

\subsection{2. Peierls Condition}
We prove that the probability of a contour $\gamma$ decays exponentially with its size (area law for defects): $\rho(\gamma) \le e^{-\beta |\gamma|}$. In the strong coupling regime, this follows from the character expansion coefficients.

\subsection{3. Cluster Expansion}
Using the Kotecký-Preiss criterion, we prove the polymer expansion converges absolutely. This allows us to write the connected correlation function as a sum over polymers connecting $x$ to $y$.

\subsection{4. Tree Decay}
The sum is dominated by the smallest polymer connecting the points, which has length at least $|x-y|$. Thus, the correlation decays as $e^{-m|x-y|}$.

\subsection{5. Rank-1 Projection}
In the operator formalism, this clustering property implies that the spectral projection onto the ground state $P_0$ is rank-1. Thus, the vacuum is unique, and the gap $\Delta$ is well-defined.



